\documentclass[a4paper,12pt]{article}

\usepackage{worksheet}
\usepackage{tikz}
\usepackage{gensymb}
\usepackage{tabularx}
\usetikzlibrary{arrows.meta}
\usetikzlibrary{decorations.markings}

\newcounter{QCounter}
\newcounter{SQCounter}[QCounter]
\newcounter{ACounter}
\newcounter{SACounter}[ACounter]

\def\qcnt{\stepcounter{QCounter}\arabic{QCounter}}
\def\sqcnt{\stepcounter{SQCounter}\alph{SQCounter}}

\def\acnt{\stepcounter{ACounter}\arabic{ACounter}}
\def\sacnt{\stepcounter{SACounter}\alph{SACounter}}

\setlength{\tabcolsep}{4mm}

\begin{document}

\section*{Standard 2 Trigonometry}

\subsection*{Question Bank}

\subsubsection*{Question \qcnt}

Victoria wants to help her friend Irene compare compound interest 
options using a table that shows the future value interest factors 
of a compound interest investment of \$1.

\begin{figure}[h]
\centering
\begin{tabular}{|c|c|c|c|c|}
\hline
& \multicolumn{4}{c|}{\emph{Rate}} \\
\hline
\emph{Period} & \emph{1\%} & \emph{2.5\%} & \emph{4\%} & \emph{5\%} \\
\hline
5 & 1.0510 & 1.1314 & 1.2167 & 1.2763 \\
\hline
10 & 1.1046 &  & 1.4802 & 1.6289 \\
\hline
20 & 1.2202 & 1.6386 & 2.1911 & 2.6533 \\
\hline
\end{tabular}
\end{figure}

\sqcnt) Victoria forgot to fill one of the entries in her table 
before giving it to Irene, what factor should go in the blank 
cell in the table?

\sqcnt) Victoria is going to invest \$4000 in an account that pays 
5\% p.a.~compounding every 6 months. Irene is going to invest \$4150
in an account that pays 4\% p.a.~compounding quarterly. Who will 
have more money in their account at the end of 5 years?

\subsubsection*{Question \qcnt}

Daphne hopes to have \$30000 in a savings account in 5 years. The 
account pays 4\% interest per annum, compounding quarterly. How 
much should she place into the account now to have the required 
amount at the end of 5 years?

\subsubsection*{Question \qcnt}

A table showing the future value interest 
factors of an annuity of \$1 is shown below.

\begin{figure}[h]
\centering
\begin{tabular}{|c|c|c|c|c|}
\hline
& \multicolumn{4}{c|}{\emph{Rate}} \\
\hline
\emph{Period} & \emph{0.05\%} & \emph{0.1\%} & \emph{0.2\%} & \emph{0.4\%} \\
\hline
135 & 139.6244 & 144.4596 & 154.8056 & 178.5399 \\
\hline
180 & 188.2993 & 197.1097 & 216.4071 & 262.8712 \\
\hline
270 & 288.9963 & 309.7877 & 357.5409 & 484.5857 \\
\hline
390 & 430.5035 & 476.6930 & 589.8868 & 936.0091 \\
\hline
\end{tabular}
\end{figure}

Cynthia is planning to save regularly and is comparing two different 
saving schemes. Scheme A is to save \$1200 monthly in an account 
that pays 2.4\% per annum compounding monthly. Scheme B is to save 
\$500 fortnightly in an account that pays 2.6\% per annum compounding 
fortnightly. Which scheme would have a higher future value after 15 
years?

\newpage
\subsubsection*{Question \qcnt}

A table showing the future value interest factors of an 
annuity of \$1 is shown below.

\begin{figure}[h]
\centering
\begin{tabular}{|c|c|c|c|c|}
\hline
& \multicolumn{4}{c|}{\emph{Rate}} \\
\hline
\emph{Period} & \emph{0.001} & \emph{0.002} & \emph{0.003} & \emph{0.004} \\
\hline
60 & 61.8047 & 63.6809 & 65.6316 & 67.6602 \\
\hline
120 & 127.4292 & 135.4722 & 144.1857 & 153.6320 \\
\hline
180 & 197.1097 & 216.4071 & 238.2067 & 262.8712 \\
\hline
240 & 271.0967 & 307.6500 & 350.7400 & 401.6750 \\
\hline
\end{tabular}
\end{figure}

Agatha plans to save \$600000 by making monthly payments into 
an account that pays 3.6\% interest per annum compounding monthly.

\sqcnt) If she wants to save \$1710.67 each month, how many years 
will it take her to achieve her goal?

\sqcnt) Agatha decides that she wants to save the amount in 
15 years, how much will she have to save each month to make 
it happen?

\subsubsection*{Question \qcnt}

Below is a table showing the present value factor for an 
annuity paying \$1.


\begin{figure}[h]
\centering
\begin{tabular}{|c|c|c|c|c|}
\hline
& \multicolumn{4}{c|}{\emph{Rate}} \\
\hline
\emph{Period} & \emph{0.001} & \emph{0.002} & \emph{0.003} & \emph{0.004} \\
\hline
60 & 58.2072 & 56.4866 & 54.8349 & 53.2489 \\
\hline
120 & 113.0264 & 106.5918 & 100.6492 & 95.1560 \\
\hline
180 & 164.6547 & 151.0364 & 138.9268 & 128.1370 \\
\hline
240 & 213.2778 & 190.4600 & 170.9076 & 154.0933 \\
\hline
300 & 259.0707 & 225.4297 & 197.6274 & 174.5210 \\
\hline
360 & 302.1982 & 256.4488 & 219.9517 & 190.5977 \\
\hline
\end{tabular}
\end{figure}

\sqcnt) Leona is to receive a pension payment of \$2400 per month 
from an account that pays 2.4\% interest per annum compounded 
monthly for 20 years. What would the present value of the plan 
be if she were to receive a lump sum of money today?

\sqcnt) Hannah takes a loan of \$60000 at an interest rate of 
4.8\% per annum compounded monthly. She is going to make monthly 
payments for 15 years, how much does she have to pay each month 
to pay off the loan?

\newpage
\subsubsection*{Question \qcnt}

Rebecca wants to save some money for the future, so she puts 
\$20000 into an account that pays 2.2\% interest per annum 
compounding monthly. After 4 years, the bank increases the 
interest rate to 2.8\% per annum and Rebecca decides to put 
an extra \$5000 into the account for the next 5 years. How 
much money will she have in the account at the end of 
the 9 years?

\subsubsection*{Question \qcnt}

Below is a table showing the future value factor for an 
annuity of \$1.

\begin{figure}[h]
\centering
\begin{tabular}{|c|c|c|c|c|}
\hline
& \multicolumn{4}{c|}{\emph{Rate}} \\
\hline
\emph{Period} & \emph{0.1\%} & \emph{0.2\%} & \emph{0.3\%} & \emph{0.4\%} \\
\hline
180 & 197.1097 & 216.4071 & 238.2067 & 262.8712 \\
\hline
216 & 240.9684 & 269.8354 & 303.2867 & 342.1367 \\
\hline
252 & 286.4340 & 327.2482 & 375.7771 & 433.6529 \\
\hline
\end{tabular}
\end{figure}

Below is a table showing the present value factor for an 
annuity paying \$1.

\begin{figure}[h]
\centering
\begin{tabular}{|c|c|c|c|c|}
\hline
& \multicolumn{4}{c|}{\emph{Rate}} \\
\hline
\emph{Period} & \emph{0.1\%} & \emph{0.2\%} & \emph{0.3\%} & \emph{0.4\%} \\
\hline
90 & 86.0277 & 82.2898 & 78.7706 & 75.4556 \\
\hline
120 & 113.0264 & 106.5918 & 100.6492 & 95.1560 \\
\hline
150 & 139.2275 & 129.4799 & 120.6474 & 112.6328 \\
\hline
\end{tabular}
\end{figure}

Lydia wants to have some money saved for when her child turns 18. 
When her child is born, she saves \$10000 into an account that pays 
3.6\% interest per annum compounding monthly. During this time she 
also saves \$100 per month into the account.

\sqcnt) When her child turns 18 how much money will be in the 
account?

\sqcnt) Once her child is 18, the account will pay out a fixed 
amount every month for the next 10 years. How much will each 
payment be?

\newpage
\subsubsection*{Question \qcnt}

Below is a table showing the present value factor for an 
annuity paying \$1.

\begin{figure}[h]
\centering
\begin{tabular}{|c|c|c|c|c|}
\hline
& \multicolumn{4}{c|}{\emph{Rate}} \\
\hline
\emph{Period} & \emph{0.0025} & \emph{0.003} & \emph{0.0035} & \emph{0.004} \\
\hline
60 & 55.6524 & 54.8349 & 54.0339 & 53.2489 \\
\hline
120 & 103.5618 & 100.6492 & 97.8489 & 95.1560 \\
\hline
180 & 144.8055 & 138.9268 & 133.3777 & 128.1370 \\
\hline
240 & 180.3109 & 170.9076 & 162.1874 & 154.0933 \\
\hline
300 & 210.8765 & 197.6274 & 185.5486 & 174.5210 \\
\hline
\end{tabular}
\end{figure}

Felicity takes out a loan with an interest rate of 4.2\% per 
annum compounded monthly. The payment plan is to pay \$1200 per 
month for the first 10 years and then \$900 per month for the 
next 15.

\sqcnt) What is the loan value after the first 10 years? 

\sqcnt) The original value of the loan would be the present value of 
value at 10 years, since it has been compounding, plus the value of 
the annuity being paid for the first 10 years. How much was 
the original value of the loan?

\sqcnt) An alternative method for calculating the original loan 
amount might not make as much sense, but we can check that it gives 
the same answer. Split the loan payments into a plan where we 
pay \$900 for 25 years, and also pay \$300 for the first 10 years. 
The original loan should be the sum of the present value of these 
two annuities. How much is it using this method?

\newpage
\subsection*{Answers}

\subsubsection*{Question \acnt}

\sacnt) Using the compound interest formula

$$(1+0.025)^{10}\approx1.2801.$$

\sacnt) Victoria's investment will use the table value 1.2801 
so her final value will be $1.2801\times4000 = \$5120.40$.

Irene's investment will use the table value 1.2202 so her 
final value will be $1.2202\times4150 = \$5063.82$. 

Hence Victoria will have more in her account at the end of 5 years.

\subsubsection*{Question \acnt}

Let her initial investment be the principle value $P$, the formula 
tells us

$$30000 = P\left(1+\frac{4\%}{4}\right)^{5\times4}=P\times1.01^{20}$$

Hence her initial amount would be 

$$P = 30000\div1.01^{20} = \$24586.33$$

\subsubsection*{Question \acnt}

For Scheme A, we need to use the factor 216.4071 giving a 
future value of $216.4071\times 1200 = \$259688.52$.

For Scheme B, use $476.6930\times 500 = \$238246.50$.

So Scheme A would have a higher future value.

\subsubsection*{Question \acnt}
\sacnt) Doing $600000\div1710.67=350.73876...\approx350.7400$, 
and referencing the table indicates that this would need 
240 periods. Since she is saving monthly 240 periods is 20 years.

\sacnt) 15 years is 180 monthly periods, and the table gives us 
a factor of 238.2067. $600000\div238.2067=\$2518.82$ saved each 
month.

\subsubsection*{Question \acnt}

\sacnt) The lump sum would be $2400\times 190.4600=\$457104$

\sacnt) $60000\div128.1370 = \$458.25$ per month.

\subsubsection*{Question \acnt}

For the amount after first 4 years $A_1$,
$$A_1 = 20000\left(1+\frac{0.022}{12}\right)^{4\times12} = \$21838.00$$
Then she saves an extra \$5000, $\$21838+\$5000=\$26838$. So after 
the next 5 years 
$$A = 26838\left(1+\frac{0.028}{12}\right)^{5\times12} = \$30866.01$$

\subsubsection*{Question \acnt}

\sacnt) This will be a combination of a compound interest problem and 
annuity problem. For the interest on the \$10000,
$$A_c = 10000\left(1+\frac{0.036}{12}\right)^{18\times12}=\$19098.60.$$
Then for the annuity, we need to use the table value 303.2867,
$$A_a = 100\times303.2867 = 30328.67.$$
The total amount in the account at the end of 18 years will be
$$A = A_c + A_a = \$19098.60 + \$30328.67 = \$49427.27.$$

\sacnt) The annuity present value is \$49427.27, and since it will pay 
\$$x$ for 10 years we will use the factor 100.6492. 
\begin{equation*}
\begin{split}
\$49427.27 &= x\times100.6492 \\
\$49427.27\div100.6492 &= x \\
\$491.08 &= x.
\end{split}
\end{equation*}

\newpage
\subsubsection*{Question \acnt}

\sacnt) Use the present value factor 133.3777, 
$$A_{10} = \$900\times133.3777 = \$120039.93.$$

\sacnt) The present value of the amount left after 10 years $A_{10}$ 
is given by the compound interest formula,
$$A_a = A_{10}\div1.0035^{10\times12} = \$120039.93\div1.0035^{10\times12} 
= \$78929.71.$$
The present value of the annuity being paid for the first 10 years
$$A_b = \$1200\times97.8489 = \$117418.68.$$
The original value of the loan would be the sum of these present values,
$$\$78929.71 + \$117418.68 = \$196348.39.$$

\sacnt) The annuity of \$900 for 25 years would be 
$$A_c = \$900\times 185.5486 = \$166993.74.$$
The annuity of \$300 for 10 years would be 
$$A_d = \$300\times 97.8489 = \$29354.67.$$
The original present value of the loan should be the sum of the 
two,
$$\$166993.74 + \$29354.67 = \$196348.41.$$
We can see this method also gets about the same amount, but might 
be a bit more difficult to understand.

\end{document}
